\section{Agent-Level GRPO}

\label{sec:agentgrpo}

\subsection{Tool-Use Trajectory Definition}
\label{sec:trajectories}

\begin{definition}[Tool-Use Trajectory]
\label{def:trajectory}
A tool-use trajectory is a tuple
\begin{equation}
    \tau = \bigl(x,\;\bm{\xi},\;R\bigr),
    \quad \bm{\xi} = \bigl((t_1, \bm{p}_1, o_1, s_1),\;\ldots,\;(t_K, \bm{p}_K, o_K, s_K)\bigr),
    \label{eq:trajectory}
\end{equation}
where:
\begin{itemize}[leftmargin=1.5em]
    \item $x \in \mathcal{X}$ is the natural-language task description,
    \item $K \geq 0$ is the number of tool calls made,
    \item $t_i \in \mathcal{T}$ is the name of the $i$-th tool invoked from
          tool registry $\mathcal{T}$,
    \item $\bm{p}_i \in \mathcal{P}(t_i)$ is the parameter dictionary for
          tool $t_i$,
    \item $o_i \in \mathcal{O}$ is the output returned by tool $t_i$,
    \item $s_i \in \{0, 1\}$ is the binary success indicator ($s_i = 1$ iff
          the call completed without error and produced well-formed output),
    \item $R \in [0, 1]$ is the trajectory-level composite reward defined in
          Section~\ref{sec:reward}.
\end{itemize}
\end{definition}

\begin{remark}
The length $K$ is not fixed a priori; it varies across trajectories in the
same group. This is fundamentally different from the text setting where all
rollouts have a common output space $\mathcal{V}^*$. Agent-GRPO must compare
trajectories of potentially different lengths, which motivates the
summarization approach in Stage~1.
\end{remark}

Let $\mathcal{T}$ be a finite registry of $|\mathcal{T}|$ available tools.
We partition $\mathcal{T}$ into \emph{domains} $\mathcal{D} =
\{d_1, \ldots, d_L\}$ based on functionality: $d_{\text{fs}}$ (file system),
$d_{\text{shell}}$ (shell execution), $d_{\text{search}}$ (web/code search),
$d_{\text{k8s}}$ (Kubernetes operations), and so on. Domain membership informs
the scoped retrieval of experiences (Section~\ref{sec:retrieval}).

\subsection{Reward Computation}
\label{sec:reward}

The trajectory-level reward $R$ decomposes into two observable components:

\paragraph{Tool Success Rate.}
The per-trajectory tool success rate is the fraction of tool calls that
completed without error:
\begin{equation}
    \rho_s(\tau) = \frac{1}{K} \sum_{i=1}^K s_i \in [0, 1].
    \label{eq:tool-success}
\end{equation}
For the degenerate case $K = 0$ (task completed without tool use), we set
$\rho_s = 1$.

\paragraph{User Acceptance Signal.}
The user acceptance signal $\rho_a(\tau) \in \{0, 1\}$ is determined by
observing whether the user corrects or rejects the agent's final output within
a time window $\Delta t$:
\begin{equation}
    \rho_a(\tau) =
    \begin{cases}
        1 & \text{if user does not correct output within } \Delta t = 120\,\text{s,} \\
        0 & \text{if user explicitly corrects or rejects output.}
    \end{cases}
    \label{eq:user-accept}
\end{equation}
This weak, delayed supervision signal is the only human signal required by
Agent-GRPO; no explicit reward labeling is needed.

\paragraph{Composite Reward.}
The composite trajectory reward is:
\begin{equation}
    R(\tau) = w_1\,\rho_s(\tau) + w_2\,\rho_a(\tau),
    \label{eq:composite-reward}
\end{equation}
with default weights $w_1 = 0.4$, $w_2 = 0.6$ (user acceptance is weighted
more heavily as it reflects task-level quality, not just mechanical correctness).

\begin{remark}
The weights $w_1$ and $w_2$ can be tuned per deployment. In automated
pipelines without a human in the loop, one may set $w_2 = 0$ and rely
entirely on $\rho_s$. In interactive settings, $w_2 > 0$ captures the
user's implicit preference signal.
\end{remark}

\paragraph{Step-Level Rewards.}
For Stage-1 summarization (Section~\ref{sec:stage1}), we also compute step-level
rewards $r_i = s_i$ for each tool call $i$, which enable the extractor to
identify exactly which calls were beneficial or harmful.

\subsection{Group-Relative Advantage for Agent Trajectories}

Given a group of $G$ trajectories $\{\tau^{(j)}\}_{j=1}^G$ on the same task
$x$, the group-relative advantage of trajectory $\tau^{(i)}$ is defined
analogously to Equation~\eqref{eq:grpo-advantage}:
\begin{equation}
    A^{(i)} = \frac{R(\tau^{(i)}) - \mu_G}{\sigma_G + \epsilon},
    \quad \mu_G = \frac{1}{G}\sum_{j=1}^G R(\tau^{(j)}),
    \quad \sigma_G = \sqrt{\frac{1}{G-1}\sum_{j=1}^G \bigl(R(\tau^{(j)}) - \mu_G\bigr)^2}.
    \label{eq:agent-advantage}
\end{equation}
Trajectories with $A^{(i)} > 0$ are identified as ``positive examples'' for
Stage-2 extraction; those with $A^{(i)} < 0$ are ``negative examples.'' The
group is required to contain at least one positive and one negative example for
meaningful advantage extraction; groups where all trajectories achieve the same
reward are discarded.

\subsection{Stage 1: Tool Trajectory Summarization}
\label{sec:stage1}

Stage~1 processes each trajectory $\tau^{(i)}$ independently, producing a
structured summary $\hat{\tau}^{(i)}$ suitable for group comparison.

\begin{definition}[Trajectory Summary]
A trajectory summary $\hat{\tau}$ is a natural-language description constrained
to at most $W_{\max} = 64$ words that identifies: (a) the task type and domain,
(b) the primary tool selection strategy employed, (c) the most significant
parameter choice that affected outcome, and (d) the root cause of any failures.
\end{definition}

The summary is produced by the following LLM call:

\begin{quote}
\textit{Prompt template for Stage 1:}
\texttt{You are analyzing a tool-use trajectory for task: \{x\}. The
trajectory made \{K\} tool calls with overall reward \{R\}. The step-by-step
calls were: \{$\xi$\}. Summarize in at most 64 words: (1) what tool selection
strategy was used, (2) what parameter choices were made for key tools, (3) what
failures occurred and why, (4) what the outcome was. Focus on patterns
generalizable to similar future tasks.}
\end{quote}

The summary $\hat{\tau}^{(i)}$ retains the advantage sign: we annotate it as
$\hat{\tau}^{(i)+}$ if $A^{(i)} > 0$ and $\hat{\tau}^{(i)-}$ if $A^{(i)} < 0$.

\subsection{Stage 2: Group Advantage Extraction}
\label{sec:stage2}

Stage~2 operates on the full group of summaries $\{\hat{\tau}^{(j)}\}_{j=1}^G$
and produces a set of \emph{operations} $\Omega^{(g)}$ on the experience library.

\begin{definition}[Experience Library Operation]
\label{def:operation}
An operation $\omega \in \Omega$ is one of:
\begin{itemize}[leftmargin=1.5em]
    \item $\textsc{add}(e_{\text{new}})$: insert a new experience $e_{\text{new}}$ into $\mathcal{E}$,
    \item $\textsc{modify}(e_{\text{id}}, e_{\text{new}})$: replace experience with identifier $e_{\text{id}}$,
    \item $\textsc{delete}(e_{\text{id}})$: remove experience $e_{\text{id}}$ from $\mathcal{E}$,
    \item $\textsc{merge}(e_{\text{id}_1}, e_{\text{id}_2}, e_{\text{new}})$: replace two near-duplicate
          experiences with a single consolidated entry.
\end{itemize}
\end{definition}

The Stage-2 LLM prompt compares positive and negative group summaries:

\begin{quote}
\textit{Prompt template for Stage 2:}
\texttt{You are extracting generalizable tool-use lessons. Below are
SUCCESSFUL trajectories (high reward): \{$\hat{\tau}^{(i)+}$\}. And FAILED
trajectories (low reward): \{$\hat{\tau}^{(i)-}$\}. Current experience library:
\{$\mathcal{E}_{\text{relevant}}$\}. Produce operations (add/modify/delete/merge)
that update the library. Each new/modified experience must: (a) be $\leq$32
words, (b) be actionable and specific, (c) explain WHEN and HOW to apply.
Format: OPERATION | experience\_id (if modify/delete/merge) | text.}
\end{quote}

The output $\Omega^{(g)}$ from Stage~2 contains all proposed library
modifications arising from task group $g$.

\subsection{Stage 3: Batch Consolidation}
\label{sec:stage3}

Stage~3 merges the operation sets $\{\Omega^{(g)}\}_{g=1}^{N_g}$ from all
$N_g$ task groups processed in the current session, producing a final
consolidated operation set $\Omega^*$ that is applied atomically to $\mathcal{E}$.

Consolidation proceeds in three passes:

\begin{enumerate}[leftmargin=1.5em]
    \item \textbf{Conflict resolution:} If two groups issue conflicting
          \textsc{modify} or \textsc{delete} operations on the same $e_{\text{id}}$,
          the operation from the group with higher mean group advantage
          $\bar{A}^{(g)} = \frac{1}{G}\sum_i |A^{(i)}_g|$ takes precedence.
    \item \textbf{Deduplication:} \textsc{add} operations whose proposed text
          has cosine similarity above $\theta_{\text{dup}} = 0.85$ with an
          existing library entry are converted to \textsc{modify} operations.
    \item \textbf{Length enforcement:} Any experience text exceeding 32 words
          is compressed via an LLM call: \texttt{Summarize in $\leq$32 words,
          preserving the actionable core: \{text\}}.
\end{enumerate}

\subsection{The AGENT-GRPO Algorithm}
\label{sec:algorithm}

Algorithm~\ref{alg:agent-grpo} presents the complete Agent-GRPO procedure.

\begin{algorithm}[H]
\caption{AGENT-GRPO: Agent-Level Training-Free GRPO}
\label{alg:agent-grpo}
\begin{algorithmic}[1]
\Require Session trajectories $\mathcal{S} = \{(\tau^{(j)}_g)\}_{g,j}$,
         experience library $\mathcal{E}$,
         group size $G$, consolidation threshold $\theta_{\text{dup}}$
\Ensure Updated experience library $\mathcal{E}'$

\State $\Omega \gets \emptyset$ \Comment{Accumulated operations}
\For{each task group $g$ in $\mathcal{S}$}
    \State $\{\tau^{(j)}\}_{j=1}^G \gets$ trajectories for group $g$
    \State \textbf{// Compute rewards and advantages}
    \For{$j = 1$ to $G$}
        \State $R^{(j)} \gets w_1\,\rho_s(\tau^{(j)}) + w_2\,\rho_a(\tau^{(j)})$
    \EndFor
    \State $\mu_G, \sigma_G \gets \text{GroupStats}(\{R^{(j)}\})$
    \If{$\sigma_G < \epsilon$}
        \State \textbf{continue} \Comment{No variance: group provides no learning signal}
    \EndIf
    \For{$j = 1$ to $G$}
        \State $A^{(j)} \gets (R^{(j)} - \mu_G) / (\sigma_G + \epsilon)$
    \EndFor

    \State \textbf{// Stage 1: Trajectory Summarization}
    \For{$j = 1$ to $G$}
        \State $\hat{\tau}^{(j)} \gets \text{Summarize}(\tau^{(j)},\, \text{sign}(A^{(j)}))$
    \EndFor

    \State \textbf{// Stage 2: Group Advantage Extraction}
    \State $\mathcal{E}_{\text{rel}} \gets \text{Retrieve}(\mathcal{E},\, x_g,\, K_{\text{stage2}}=5)$
    \State $\Omega^{(g)} \gets \text{ExtractAdvantages}(\{\hat{\tau}^{(j)}\},\, \mathcal{E}_{\text{rel}})$
    \State $\Omega \gets \Omega \cup \Omega^{(g)}$
\EndFor

\State \textbf{// Stage 3: Batch Consolidation}
\State $\Omega^* \gets \text{Consolidate}(\Omega,\, \theta_{\text{dup}})$

\State \textbf{// Apply operations}
\State $\mathcal{E}' \gets \text{Apply}(\mathcal{E},\, \Omega^*)$

\State \Return $\mathcal{E}'$
\end{algorithmic}
\end{algorithm}

\paragraph{Complexity.}
Let $N_g$ be the number of task groups, $G$ the group size, and $L_{\text{LLM}}$
the cost of a single LLM extraction call. Stage~1 requires $N_g G$ LLM calls.
Stage~2 requires $N_g$ calls. Stage~3 requires $O(|\Omega|)$ comparisons plus
at most $|\Omega|$ LLM compression calls. Total cost: $O(N_g(G + 1))$ LLM
calls, typically 50--200 calls per session at approximately \$0.05--\$0.10
each, yielding a per-session adaptation cost of \$2.50--\$20.

% ============================================================================
% 4. SEMANTIC MEMORY MANAGER
% ============================================================================
